量化交易在金融市场中的主导地位日益增强，交易策略已构成一个复杂的、相互依赖的"策略生态系统"。在此背景下，策略超额收益（Alpha）是否会因策略间的竞争、模仿和替代而系统性消失，成为兼具理论价值和实践意义的核心问题。本研究整合演化金融学、适应性市场假说和基于Agent的计算经济学三个理论框架，构建演化人工股票市场模型（EVO-ASM），系统性地研究策略生态系统中的Alpha消失机制。

模型包含500个异质性Agent、6维信号空间、Walrasian市场出清机制以及含有淘汰-复制-突变操作的演化引擎。通过四组递进实验——无演化基线（E1）、资本扩张（E2）、复制机制（E3）、复制+突变完整演化（E4）——逐步剥离Alpha衰减的因果机制，共计完成数百次独立模拟运行。

实验结果表明：（1）Alpha衰减需要演化机制驱动，在静态策略分布下Alpha不会系统性消失；（2）财富集中单独即可驱动Alpha部分衰减，拥挤度每上升1个百分点，Alpha下降约0.3-0.8个基点；（3）策略复制使策略多样性丧失速度加快2-3倍；（4）突变是维持策略生态多样性的关键，适度突变可在竞争与创新之间建立动态平衡；（5）策略生态系统存在临界相变，系统可能从多样化状态突然跃迁至同质化状态，此发现对金融稳定具有重要预警意义。本研究将Alpha衰减从实证残差重新概念化为演化涌现属性，为理解量化时代的市场效率动态提供了新的理论框架和分析工具。
